\section*{第 \ref{ch:g9:solids} 章 --- 立体图形与体积}

\begin{solution}{exo:g9:solids:1}
长方体：$V = 4 \times 5 \times 12 = 240$。

圆柱：$V = \pi r^2 h = \pi \times 9 \times 7 = 63\pi \approx 198$。
\end{solution}

\begin{solution}{exo:g9:solids:2}
圆锥：$V = \frac13 \pi r^2 h = \frac13 \pi \times 25 \times 12 =
100\pi \approx 314$。

棱锥：底面积 $B = 8 \times 3 = 24$，所以
$V = \frac13 \times 24 \times 10 = 80$。
\end{solution}

\begin{solution}{exo:g9:solids:3}
体积：$V = \frac43 \pi \times 6^3 = \frac43 \pi \times 216 = 288\pi$。

表面积：$4\pi \times 6^2 = 144\pi$。
\end{solution}

\begin{solution}{exo:g9:solids:4}
截面半径：$\sqrt{r^2 - d^2} = \sqrt{100 - 64} = \sqrt{36}
= 6$。
\end{solution}

\begin{solution}{exo:g9:solids:5}
从整个圆锥到小圆锥的相似比：
$k = \frac{9}{12} = \frac34$。截面半径：$8 \times \frac34 = 6$。

整个圆锥的体积：$\frac13\pi \times 64 \times 12 = 256\pi$。
小圆锥的体积：$256\pi \times \left(\frac34\right)^3 =
256\pi \times \frac{27}{64} = 108\pi \approx 339$。
\end{solution}

\begin{solution}{exo:g9:solids:6}
球的体积：$\frac43\pi \times 2^3 = \frac{32\pi}{3}$ cm$^3$。
这部分体积铺在底面积 $\pi \times 3^2 = 9\pi$ cm$^2$ 上，
因此水位上升
\[
\frac{32\pi/3}{9\pi} = \frac{32}{27} \approx 1.2 \text{ cm}
\]
（约 $12$ mm）。
\end{solution}

\begin{solution}{exo:g9:solids:7}
小模具是大模具按比例 $k = \frac36 =
\frac12$，因此其体积是大模具的 $\left(\frac12\right)^3 = \frac18$，
配方可以装满 $8$ 个小模具。
\end{solution}

\begin{solution}{exo:g9:solids:8}
$V = \frac13 \times 230^2 \times 147 = \frac13 \times 52\,900 \times 147
= 52\,900 \times 49 = 2\,592\,100$ m$^3$ ——约为 $2.6$ 百万
立方米。
\end{solution}

\begin{solution}{exo:g9:solids:9}
\emph{1.} 水形成一个按 $k = \frac12$ 缩放的圆锥（尖端朝下，
高为一半），因此其体积为漏斗的 $\left(\frac12\right)^3 = \frac18$
：只有八分之一，远少于一半！

\emph{2.} 水深为 $d$ 时形成一个按
$k = \frac{d}{12}$ 缩放的圆锥，其体积为漏斗的 $k^3$ 倍。体积为一半意味着
$k^3 = \frac12$，即
$\left(\frac{d}{12}\right)^3 = \frac12$。于是
$\frac{d}{12} = \sqrt[3]{0.5} \approx 0.7937$，因此
$d \approx 9.5$ cm：后一半体积只占据顶部
$2.5$ cm——锥体较宽的部分容纳了大部分的水。
\end{solution}

\begin{solution}{pb:g9:solids:1}
\textbf{1.} $V_{\text{圆柱}} = \pi r^2 \times 2r = 2\pi r^3$。

\textbf{2.} $\dfrac{V_{\text{球}}}{V_{\text{圆柱}}}
= \dfrac{\frac43 \pi r^3}{2 \pi r^3} = \dfrac23$。$\pi$ 与
$r^3$ 约去：这一比例是\emph{普遍的}——
对弹珠和行星都成立。这是形状之间的关系，而非
数字之间的关系：正是值得刻在
石头上的那类真理。

\textbf{3.} 球：$4\pi r^2$。圆柱：侧面积
$2\pi r \times 2r = 4\pi r^2$，加上两个底面
$2 \times \pi r^2$，总计 $6\pi r^2$。比值：
$\frac{4\pi r^2}{6\pi r^2} = \frac23$——同样的三分之二，
这次针对表面积。（还有一个附带结论：球的表面积恰好等于
圆柱的侧面积。）

\textbf{4.} 剩余体积：$2\pi r^3 - \frac43\pi r^3 =
\frac23 \pi r^3$。两个半径为 $r$、高为 $r$ 的圆锥：
$2 \times \frac13 \pi r^2 \times r = \frac23 \pi r^3$。两者相等。

\textbf{5.} 球：$\frac43 \pi \times 12^3 \approx 7\,200$
cm$^3$；盒子：$2\pi \times 12^3 \approx 10\,900$ cm$^3$。空的部分：
占盒子的三分之一，约为 $33\,\%$——
由第 2 题保证，与球的大小无关。

\textbf{6.} 半球：
$\frac12 \times \frac43 \pi \times 10^3 = \frac23 \pi \times
1\,000 \approx 2\,094$ cm$^3 \approx 2.1$ L。

\textbf{7.} 圆锥：$\frac13 \pi \times 10^2 \times 10 =
\frac{1\,000\pi}{3} \approx 1\,047$ cm$^3$。半球：
$\frac{2\,000\pi}{3} \approx 2\,094$ cm$^3$。圆柱：
$1\,000\pi \approx 3\,142$ cm$^3$。比值：
$\frac{1000\pi}{3} : \frac{2000\pi}{3} : \frac{3000\pi}{3}
= 1 : 2 : 3$，恰好成立。

\textbf{8.} $\frac43 \pi \times 27 = \pi \times 9 \times h$
得 $h = \frac43 \times 3 = 4$ cm，恰好成立。

\textbf{9.} 半径：$\frac{6\,371}{1\,737} \approx 3.67$。
体积按立方比例变化（见\cref{thm:g9:solids:scaling}）：
$3.67^3 \approx 49$。地球中大约可以容纳五十个月球。

\textbf{10.} 表面积按平方比例变化：$3.67^2 \approx
13.5$。将气球的半径加倍，其体积变为原来的
$2^3 = 8$ 倍，而表面积仅变为原来的 $2^2 = 4$ 倍：
随着尺寸增大，体积（重量、容量、产生的热量）增长得快于表面（外皮、
材料、散失的热量） --- 平方--立方定律。

\textbf{11.} 重量：$\times 10^3 = 1\,000$。骨骼
截面：$\times 10^2 = 100$。压强 --- 骨骼单位
面积上的重量：$\times \frac{1000}{100} = 10$。按人体压强建造的
骨骼需承受十倍的压强：巨人的骨架
会在自身重量下折断。巨人在几何上是
不可能存在的；大型动物需要不成比例地粗壮的骨骼
（比较大象的腿与瞪羚的腿）。

\textbf{12.} 重量：$\div 100^3 = 10^6$。力量（肌肉
截面）：$\div 100^2 = 10^4$。力量与重量之比：
$\times \frac{10^6}{10^4} = 100$。能举起自身五十倍重量的蚂蚁
并非超级运动员 --- 它只是体型小；
若把人缩小到蚂蚁的体型，也能做同样的事。

\textbf{13.} $\dfrac{S}{V} = \dfrac{4\pi r^2}{\frac43\pi r^3}
= \dfrac{3}{r}$：当 $r = 1$ 时比值为 $3$，当 $r = 2$ 时
比值为 $1.5$ --- 尺寸减半，使单位产热体积
所对应的可用表面积加倍。小型生物散热很快：
老鼠必须不停地进食，熊可以冬眠，而
婴儿 --- 小球体，$\frac Sr$ 较大 --- 需要戴帽子。

\textbf{14.} $\left(\frac54\right)^3 = \frac{125}{64} \approx
1.95$：同样价格下，大橙子所含果肉几乎是\emph{两倍}。
因此选购时应看半径。

\textbf{15.} 石刻上写道：球的体积与圆柱体积之比 $= \frac23$，
\emph{并且}总表面积也是其 $\frac23$ ---
两个精确、普遍的比例，把最圆的立体与最简单的立体联系起来。
西塞罗的回答是：因为定理是唯一一种
既不会被军队毁坏、也不会被岁月侵蚀的纪念碑，
阿基米德的器械随叙拉古一同被焚毁，
但这两个三分之二仍然恰好是三分之二。
\end{solution}
